\begin{proof}[Proof of \cref{obj:lem:observable-marked-reduction}]
We establish the observable identities, the lower bounds, the functional representation, and then the derivative bounds.

\begin{enumerate}
\item For the target coordinate, \cref{obj:def:marked-density-vector} gives \(F_0=f_{\mathrm T}\). By \cref{obj:ass:target-density-holder}, this density is continuous on \([0,1]\), hence measurable there. For every measurable \(B\subseteq[0,1]\), \cref{obj:ass:target-density-bounds} therefore yields
\[
P_{\mathrm T}(B)
=\int_B\max\{f_{\mathrm T}(x),0\}\,dx
=\int_B f_{\mathrm T}(x)\,dx
=\int_B F_0(x)\,dx.
\]
The second equality uses \(f_{\mathrm T}(x)\ge c_f>0\) on \(B\).
% lean: CausalSmith.Stat.TransportCaceRoughnuisanceLength.observable_marked_reduction

\item For the source coordinates, define the encouragement-arm probabilities by
\[
\pi_z(x)=
\begin{cases}
1-e(x),&z=0,\\
e(x),&z=1,
\end{cases}
\qquad x\in[0,1],\quad z\in\{0,1\}.
\]
By the definition of the source propensity,
\[
P_{\mathrm S}(Z=z\mid X)=\pi_z(X)
\qquad P_{\mathrm S}\text{-almost surely}.
\]
Conditioning on \(X\) under the observed source law therefore gives
\[
P_{\mathrm S}(X\in B,Z=z)
=
\mathbb E_{P_{\mathrm S}}\!\left[
\mathbf 1_{\{X\in B\}}\pi_z(X)
\right]
=
\int_B f_{\mathrm S}(x)\pi_z(x)\,dx.
\]
Here the last equality uses the source covariate density identity in
\cref{obj:ass:source-density-bounds}; measurability follows from the stated regularity assumptions. Integrating an arm indicator also gives
\[
\int_{\{X\in B\}}\mathbf 1_{\{Z=z\}}\,dP_{\mathrm S}
=P_{\mathrm S}(X\in B,Z=z).
\]
Since \(Z\) is binary, \(\mathbf 1_{\{Z=0\}}=1-Z\) and
\(\mathbf 1_{\{Z=1\}}=Z\). Thus the two assignment coordinates satisfy
\[
\int_B F_1(x)\,dx
=\int_B f_{\mathrm S}(x)(1-e(x))\,dx
=\int_{\{X\in B\}}(1-Z)\,dP_{\mathrm S},
\]
\[
\int_B F_2(x)\,dx
=\int_B f_{\mathrm S}(x)e(x)\,dx
=\int_{\{X\in B\}}Z\,dP_{\mathrm S}.
\]

For either response \(A\in\{Y,D\}\), restriction by the same indicator and the source arm-moment identity in \cref{obj:ass:arm-means-holder} give
\[
\begin{aligned}
\int_{\{X\in B\}}\mathbf 1_{\{Z=z\}}A\,dP_{\mathrm S}
&=\int_{\{X\in B,\;Z=z\}}A\,dP_{\mathrm S}\\
&=\int_B f_{\mathrm S}(x)\pi_z(x)m_{Az}(x)\,dx
=\int_B r_{Az}(x)\,dx.
\end{aligned}
\]
Applying this identity first to \(Y\) in arms \(0,1\), and then to \(D\) in arms \(0,1\), yields respectively
\[
\begin{aligned}
\int_B F_3(x)\,dx&=\int_{\{X\in B\}}(1-Z)Y\,dP_{\mathrm S},&
\int_B F_4(x)\,dx&=\int_{\{X\in B\}}ZY\,dP_{\mathrm S},\\
\int_B F_5(x)\,dx&=\int_{\{X\in B\}}(1-Z)D\,dP_{\mathrm S},&
\int_B F_6(x)\,dx&=\int_{\{X\in B\}}ZD\,dP_{\mathrm S}.
\end{aligned}
\]
These arguments apply also when \(B\) is empty.
% lean: CausalSmith.Stat.TransportCaceRoughnuisanceLength.source_assignment_density
% lean: CausalSmith.Stat.TransportCaceRoughnuisanceLength.source_assignment_mark_integral
% lean: CausalSmith.Stat.TransportCaceRoughnuisanceLength.source_marked_arm_density

\item Fix \(x\in[0,1]\). By
\cref{obj:ass:source-density-bounds,obj:ass:propensity-overlap},
\(f_{\mathrm S}(x)\ge c_f>0\) and both \(e(x)\) and \(1-e(x)\) are at least \(1/4\). Consequently,
\[
\frac{c_f}{4}
=c_f\frac14
\le f_{\mathrm S}(x)\frac14
\le f_{\mathrm S}(x)(1-e(x))
=q_0(x)=F_1(x),
\]
\[
\frac{c_f}{4}
=c_f\frac14
\le f_{\mathrm S}(x)\frac14
\le f_{\mathrm S}(x)e(x)
=q_1(x)=F_2(x).
\]
% lean: CausalSmith.Stat.TransportCaceRoughnuisanceLength.markedDensityVector_assignment_lower_bounds

\item The preceding bounds make both arm densities strictly positive. The identities in \cref{obj:def:marked-density-vector} therefore permit cancellation:
\[
\frac{r_{Az}(x)}{q_z(x)}=m_{Az}(x),
\qquad A\in\{Y,D\},\quad z\in\{0,1\},\quad x\in[0,1].
\]
Substituting into the integrand of \cref{obj:def:smooth-functional} gives, for each response,
\[
\Phi_A(F(x))
=f_{\mathrm T}(x)\bigl(m_{A1}(x)-m_{A0}(x)\bigr).
\]
Integrating this pointwise equality proves
\[
T_A
=\int_{[0,1]}f_{\mathrm T}(x)
       \bigl(m_{A1}(x)-m_{A0}(x)\bigr)\,dx
=\int_0^1\Phi_A(F(x))\,dx.
\]
The closed-interval set integral and the displayed interval integral agree because their endpoint difference has Lebesgue measure zero.
% lean: CausalSmith.Stat.TransportCaceRoughnuisanceLength.Phi_markedDensityVector_eq_transport_integrand
% lean: CausalSmith.Stat.TransportCaceRoughnuisanceLength.transportedForm_eq_integral_Phi_markedDensityVector

\item Fix \(A\in\{Y,D\}\) and a vector \(v\) in the clipping rectangle. Define the reciprocal scale and the desired envelope by
\[
R_{\mathrm{inv}}=\frac4{c_f},
\qquad
M_{\mathrm{der}}=48(1+C_f)^2R_{\mathrm{inv}}^5.
\]
The hypotheses imply \(R_{\mathrm{inv}}>1\) and \(1+C_f>0\). To identify the response coordinates, set
\[
j_{Az}=
\begin{cases}
3+z,&A=Y,\\
5+z,&A=D,
\end{cases}
\qquad A\in\{Y,D\},\quad z\in\{0,1\}.
\]
On the open set \(\{v\in\mathbb R^7:v_1>0,\ v_2>0\}\), define the two arm contributions by
\[
H_{Az}(v)=\frac{v_0v_{j_{Az}}}{v_{1+z}},
\qquad
\Phi_A(v)=H_{A1}(v)-H_{A0}(v).
\]
The clipping rectangle lies in this open set.

For order zero, the rectangle supplies
\[
0\le v_0\le C_f,\qquad
0\le v_{j_{Az}}\le\frac{3C_f}{4},\qquad
v_{1+z}\ge\frac{c_f}{4}>0.
\]
In particular,
\[
v_0v_{j_{Az}}\le(1+C_f)^2,
\qquad
1=\frac{c_f}{4}R_{\mathrm{inv}}
\le v_{1+z}R_{\mathrm{inv}},
\]
and hence
\[
0\le H_{Az}(v)\le(1+C_f)^2R_{\mathrm{inv}}.
\]
Since
\[
R_{\mathrm{inv}}
=R_{\mathrm{inv}}\cdot1
\le R_{\mathrm{inv}}R_{\mathrm{inv}}^4
=R_{\mathrm{inv}}^5,
\]
we obtain
\[
|\Phi_A(v)|
\le 2(1+C_f)^2R_{\mathrm{inv}}
\le48(1+C_f)^2R_{\mathrm{inv}}^5
=M_{\mathrm{der}}.
\]
This treats the empty list of coordinate indices at order zero.
% lean: CausalSmith.Stat.TransportCaceRoughnuisanceLength.Phi_value_abs_le_fourthDerivativeEnvelope

\item For the positive orders, fix arbitrary coordinate indices
\(i_1,\ldots,i_4\in\{0,\ldots,6\}\), allowing repetitions. Every shorter list can be extended to such a list. Define its coordinate directions by
\[
\eta^{(\ell)}_j=
\begin{cases}
1,&j=i_\ell,\\
0,&j\ne i_\ell,
\end{cases}
\qquad 1\le\ell\le4,\quad 0\le j\le6.
\]
Thus \(|\eta^{(\ell)}_j|\le1\). For a fixed arm \(A,z\), define
\[
\mathsf U_\ell
=\eta^{(\ell)}_0v_{j_{Az}}
 +v_0\eta^{(\ell)}_{j_{Az}},
\qquad 1\le\ell\le4.
\]
The rectangle gives
\[
|v_0|\le C_f,\qquad |v_{j_{Az}}|\le C_f,\qquad
0<v_{1+z},\qquad
|v_{1+z}^{-1}|
=v_{1+z}^{-1}
\le\left(\frac{c_f}{4}\right)^{-1}
=R_{\mathrm{inv}}.
\]
The product and reciprocal rules yield
\[
\partial_{i_1}H_{Az}(v)
=\mathsf U_1v_{1+z}^{-1}
-v_0v_{j_{Az}}\eta^{(1)}_{1+z}v_{1+z}^{-2}.
\]
Moreover,
\[
|\mathsf U_1|
\le|\eta^{(1)}_0|\,|v_{j_{Az}}|
   +|v_0|\,|\eta^{(1)}_{j_{Az}}|
\le2C_f.
\]
The two terms therefore obey
\[
|\mathsf U_1v_{1+z}^{-1}|\le2C_fR_{\mathrm{inv}},
\qquad
|v_0v_{j_{Az}}\eta^{(1)}_{1+z}v_{1+z}^{-2}|
\le C_f^2R_{\mathrm{inv}}^2.
\]
Using
\[
R_{\mathrm{inv}}\le R_{\mathrm{inv}}^2,\qquad
C_f\le(1+C_f)^2,\qquad
C_f^2\le(1+C_f)^2,
\]
we obtain the one-arm bound
\[
|\partial_{i_1}H_{Az}(v)|
\le2C_fR_{\mathrm{inv}}+C_f^2R_{\mathrm{inv}}^2
\le3(1+C_f)^2R_{\mathrm{inv}}^2.
\]
Both arm functions are smooth on the stated open set, so differentiation of their difference and \(R_{\mathrm{inv}}^2\le R_{\mathrm{inv}}^5\) give
\[
|\partial_{i_1}\Phi_A(v)|
\le|\partial_{i_1}H_{A1}(v)|
  +|\partial_{i_1}H_{A0}(v)|
\le6(1+C_f)^2R_{\mathrm{inv}}^2
\le M_{\mathrm{der}}.
\]
% lean: CausalSmith.Stat.TransportCaceRoughnuisanceLength.iteratedFDeriv_one_Phi_apply
% lean: CausalSmith.Stat.TransportCaceRoughnuisanceLength.marked_arm_first_derivative_abs_le
% lean: CausalSmith.Stat.TransportCaceRoughnuisanceLength.Phi_first_derivative_abs_le_fourthDerivativeEnvelope

\item For order two, define the paired direction coefficients by
\[
\mathsf V_{\ell,\ell'}
=\eta^{(\ell)}_0\eta^{(\ell')}_{j_{Az}}
 +\eta^{(\ell')}_0\eta^{(\ell)}_{j_{Az}},
\qquad 1\le\ell,\ell'\le4.
\]
Their bounds, and those of the linear coefficients, are
\[
|\mathsf V_{\ell,\ell'}|
\le1\cdot1+1\cdot1=2,
\qquad
|\mathsf U_\ell|\le2C_f.
\]
As functions of \(v\), these coefficients satisfy
\[
\partial_{i_{\ell'}}\mathsf U_\ell
=\mathsf V_{\ell,\ell'},
\qquad
\partial_{i_{\ell'}}\mathsf V_{\ell,\nu}=0,
\qquad
\partial_{i_\ell}v_{1+z}^{-p}
=-p\,\eta^{(\ell)}_{1+z}v_{1+z}^{-(p+1)}
\quad(p\ge1).
\]
Differentiating the first-order expression thus gives
\[
\begin{aligned}
\partial_{i_1}\partial_{i_2}H_{Az}(v)
={}&\mathsf V_{1,2}v_{1+z}^{-1}
-\mathsf U_1\eta^{(2)}_{1+z}v_{1+z}^{-2}
-\mathsf U_2\eta^{(1)}_{1+z}v_{1+z}^{-2}\\
&+2v_0v_{j_{Az}}\eta^{(1)}_{1+z}
                      \eta^{(2)}_{1+z}v_{1+z}^{-3}.
\end{aligned}
\]
The four terms have absolute values bounded respectively by
\[
2R_{\mathrm{inv}},\qquad
2C_fR_{\mathrm{inv}}^2,\qquad
2C_fR_{\mathrm{inv}}^2,\qquad
2C_f^2R_{\mathrm{inv}}^3.
\]
Consequently,
\[
\begin{aligned}
|\partial_{i_1}\partial_{i_2}H_{Az}(v)|
&\le2R_{\mathrm{inv}}
    +4C_fR_{\mathrm{inv}}^2
    +2C_f^2R_{\mathrm{inv}}^3\\
&\le(2+4+2)(1+C_f)^2R_{\mathrm{inv}}^5\\
&\le24(1+C_f)^2R_{\mathrm{inv}}^5.
\end{aligned}
\]
Here each reciprocal power is at most \(R_{\mathrm{inv}}^5\), and each of
\(1,C_f,C_f^2\) is at most \((1+C_f)^2\). Subtracting the two arm derivatives gives
\[
|\partial_{i_1}\partial_{i_2}\Phi_A(v)|
\le|\partial_{i_1}\partial_{i_2}H_{A1}(v)|
  +|\partial_{i_1}\partial_{i_2}H_{A0}(v)|
\le M_{\mathrm{der}}.
\]
% lean: CausalSmith.Stat.TransportCaceRoughnuisanceLength.cubicRatioD2_basis_abs_le
% lean: CausalSmith.Stat.TransportCaceRoughnuisanceLength.dPhi2_abs_le

\item For order three, the rectangle also supplies the nonnegative numerator bounds
\[
0\le v_0\le C_f,\qquad
0\le v_{j_{Az}}\le C_f,
\qquad
0<v_{1+z},\qquad
0<v_{1+z}^{-1}\le R_{\mathrm{inv}}.
\]
Differentiating the second-order expression by the same product and reciprocal rules yields
\[
\begin{aligned}
\partial_{i_1}\partial_{i_2}\partial_{i_3}H_{Az}(v)
={}&-\mathsf V_{1,2}\eta^{(3)}_{1+z}v_{1+z}^{-2}\\
&-\bigl(\mathsf V_{1,3}\eta^{(2)}_{1+z}
       +\mathsf V_{2,3}\eta^{(1)}_{1+z}\bigr)v_{1+z}^{-2}\\
&+2\mathsf U_1\eta^{(2)}_{1+z}\eta^{(3)}_{1+z}v_{1+z}^{-3}\\
&+2\mathsf U_2\eta^{(1)}_{1+z}\eta^{(3)}_{1+z}v_{1+z}^{-3}\\
&+2\mathsf U_3\eta^{(1)}_{1+z}\eta^{(2)}_{1+z}v_{1+z}^{-3}\\
&-6v_0v_{j_{Az}}\eta^{(1)}_{1+z}
                    \eta^{(2)}_{1+z}\eta^{(3)}_{1+z}v_{1+z}^{-4}.
\end{aligned}
\]
The first paired coefficient has absolute value at most \(2\), and the grouped coefficient satisfies
\[
\left|\mathsf V_{1,3}\eta^{(2)}_{1+z}
       +\mathsf V_{2,3}\eta^{(1)}_{1+z}\right|
\le2\cdot1+2\cdot1=4.
\]
Thus the six displayed terms have absolute values bounded respectively by
\[
2R_{\mathrm{inv}}^2,\qquad
4R_{\mathrm{inv}}^2,\qquad
4C_fR_{\mathrm{inv}}^3,\qquad
4C_fR_{\mathrm{inv}}^3,\qquad
4C_fR_{\mathrm{inv}}^3,\qquad
6C_f^2R_{\mathrm{inv}}^4.
\]
Adding these bounds gives
\[
\begin{aligned}
|\partial_{i_1}\partial_{i_2}\partial_{i_3}H_{Az}(v)|
&\le6R_{\mathrm{inv}}^2
    +12C_fR_{\mathrm{inv}}^3
    +6C_f^2R_{\mathrm{inv}}^4\\
&\le(6+12+6)(1+C_f)^2R_{\mathrm{inv}}^5\\
&=24(1+C_f)^2R_{\mathrm{inv}}^5.
\end{aligned}
\]
The second inequality uses \(R_{\mathrm{inv}}^2,R_{\mathrm{inv}}^3,
R_{\mathrm{inv}}^4\le R_{\mathrm{inv}}^5\) and the same numerator bounds as before. For \(A=D\), the two arm terms use the coordinate triples \((0,6,2)\) and \((0,5,1)\); for \(A=Y\), they use \((0,4,2)\) and \((0,3,1)\). In either case,
\[
|\partial_{i_1}\partial_{i_2}\partial_{i_3}\Phi_A(v)|
\le|\partial_{i_1}\partial_{i_2}\partial_{i_3}H_{A1}(v)|
  +|\partial_{i_1}\partial_{i_2}\partial_{i_3}H_{A0}(v)|
\le M_{\mathrm{der}}.
\]
% lean: CausalSmith.Stat.TransportCaceRoughnuisanceLength.cubicRatioD3_basis_abs_le
% lean: CausalSmith.Stat.TransportCaceRoughnuisanceLength.dPhi3_abs_le

\item The remaining order is four. Differentiating the third-order expression gives
\[
\begin{aligned}
\partial_{i_1}\partial_{i_2}\partial_{i_3}\partial_{i_4}H_{Az}(v)
={}&
2v_{1+z}^{-3}
\sum_{1\le\ell<\ell'\le4}
\mathsf V_{\ell,\ell'}
\prod_{\substack{1\le\nu\le4\\\nu\ne\ell,\ell'}}
\eta^{(\nu)}_{1+z}\\
&-6v_{1+z}^{-4}
\sum_{\ell=1}^4
\mathsf U_\ell
\prod_{\substack{1\le\nu\le4\\\nu\ne\ell}}
\eta^{(\nu)}_{1+z}\\
&+24v_0v_{j_{Az}}v_{1+z}^{-5}
\prod_{\nu=1}^4\eta^{(\nu)}_{1+z}.
\end{aligned}
\]
There are six pairs in the first sum. Each corresponding term satisfies
\[
\left|
2\mathsf V_{\ell,\ell'}
\prod_{\substack{1\le\nu\le4\\\nu\ne\ell,\ell'}}
\eta^{(\nu)}_{1+z}v_{1+z}^{-3}
\right|
\le2\cdot2\cdot R_{\mathrm{inv}}^3
=4R_{\mathrm{inv}}^3.
\]
Each of the four terms in the second sum satisfies
\[
\left|
6\mathsf U_\ell
\prod_{\substack{1\le\nu\le4\\\nu\ne\ell}}
\eta^{(\nu)}_{1+z}v_{1+z}^{-4}
\right|
\le6\cdot2C_fR_{\mathrm{inv}}^4
=12C_fR_{\mathrm{inv}}^4.
\]
The final term is bounded by
\[
\left|
24v_0v_{j_{Az}}v_{1+z}^{-5}
\prod_{\nu=1}^4\eta^{(\nu)}_{1+z}
\right|
\le24C_f^2R_{\mathrm{inv}}^5.
\]
Therefore,
\[
\begin{aligned}
|\partial_{i_1}\partial_{i_2}\partial_{i_3}\partial_{i_4}H_{Az}(v)|
&\le24R_{\mathrm{inv}}^3
    +48C_fR_{\mathrm{inv}}^4
    +24C_f^2R_{\mathrm{inv}}^5\\
&\le24R_{\mathrm{inv}}^5
    +48C_fR_{\mathrm{inv}}^5
    +24C_f^2R_{\mathrm{inv}}^5\\
&=24(1+C_f)^2R_{\mathrm{inv}}^5.
\end{aligned}
\]
The middle inequality uses \(R_{\mathrm{inv}}>1\). Subtracting the two arm derivatives finally gives
\[
\begin{aligned}
|\partial_{i_1}\partial_{i_2}\partial_{i_3}\partial_{i_4}\Phi_A(v)|
&\le
|\partial_{i_1}\partial_{i_2}\partial_{i_3}\partial_{i_4}H_{A1}(v)|
+
|\partial_{i_1}\partial_{i_2}\partial_{i_3}\partial_{i_4}H_{A0}(v)|\\
&\le48(1+C_f)^2R_{\mathrm{inv}}^5
=48(1+C_f)^2\left(\frac4{c_f}\right)^5.
\end{aligned}
\]
Together with the preceding cases, this proves the asserted bound for every order \(k\in\{0,\ldots,4\}\) and every choice of coordinate indices.
% lean: CausalSmith.Stat.TransportCaceRoughnuisanceLength.cubicRatioD4_basis_abs_le
% lean: CausalSmith.Stat.TransportCaceRoughnuisanceLength.dPhi4_abs_le
\end{enumerate}
\end{proof}
